\section{Introduction}

% =============================================================================
% GOAL (2-3 pages).
%
% The narrative arc to preserve:
%   I/O is a bottleneck -> IO500 measures it -> the score depends on the
%   configuration -> tuning it properly is intractable -> we make it cheap ->
%   here is what we found.
%
% Written so far: context and problem statement, up to the research question.
% Still to write: objectives, contributions, organisation of the report (the
% bullets at the end of this file). Follow the same style: continuous prose,
% one idea per paragraph, no bullet lists in the final text.
% =============================================================================

High Performance Computing (HPC) systems are used to run large-scale
scientific and engineering applications.
\replace{These applications often process very large amounts of data and use
many computing resources in parallel.}{The applications use many compute
resources in parallel and often process large amounts of data.}
\add{Therefore,} \replace{their performance}{an application's performance} does not
only depend on the computing power of the system\add{, but also on the
performance of its storage infrastructure}.
\del{The ability to access and store data efficiently is also an important
factor.}
In many HPC workloads, I/O operations \replace{can become a bottleneck and
limit the}{become the bottleneck that limits} overall application performance.

\del{This problem is becoming more important as the amount of data produced and
processed by HPC applications continues to increase.
The growth of data can be faster than the ability of storage systems to absorb
and process it.
Improving I/O performance is therefore an important research and engineering
problem in HPC.}

Parallel file systems \add{(PFS)} such as Lustre\add{~\cite{lustre}},
BeeGFS\add{~\cite{beegfs}} and GPFS\add{~\cite{gpfs}} are commonly used to
provide high-performance storage in HPC environments.
\add{In a PFS, the storage capacity is exposed as a set of logical units called
Object Storage Targets (OSTs). When a file is written, it is split into fixed-size chunks that
are sent in parallel to several targets, so that the bandwidth of many storage
devices is used at once. This process is known as striping and is illustrated
in Figure~\ref{fig:striping}. How many targets a file is spread over, called
the stripe count, and how large each chunk is, called the stripe size, are
configuration choices rather than fixed properties of the system.}
\add{Therefore,} \replace{These systems}{a PFS} exposes configuration parameters
that can affect I/O performance.
\replace{For example}{Beyond striping}, performance can \add{also} depend on
\del{the striping configuration, the number of storage targets used,} the
transfer size, the number of processes, or the placement of processes on
compute nodes.
The best values for these parameters are not necessarily the same for every
machine: a configuration that performs well on one system may perform poorly
on another one.

\begin{figure}[ht]
    \centering
    \includegraphics[width=0.62\textwidth]{figures/striping.png}
    \caption{Striping in a parallel file system. A file created on a compute
             node is split into fixed-size chunks, which are written in
             parallel to storage targets hosted by the storage servers. Here
             the stripe count is three, so the file is spread over three of
             the eight available targets.}
    \label{fig:striping}
\end{figure}


\add{To evaluate the performance of these systems, we rely on benchmarks.
A benchmark is a program designed to stress a specific part of a system under
controlled and reproducible conditions, so that the measurements it produces
can be compared across machines and configurations~\cite{hoefler}.
For I/O, two of the most widely used benchmarks are IOR~\cite{ior}, which
measures the bandwidth reached by large parallel read and write operations,
and mdtest~\cite{mdtest}, which measures the rate at which metadata operations
such as creating, opening and removing files are performed.
A different approach is taken by mini-applications, or miniapps~\cite{mantevo},
which reproduce the access pattern of a real scientific application in a
reduced and self-contained form, and therefore measure performance under a
workload closer to production use.
Each of these benchmarks characterises one aspect of the system, and none of
them alone describes the storage system as a whole.}

\replace{This}{Thus, the variety of parameters and access patterns} makes the
evaluation of I/O performance difficult\replace{.
A}{: a} benchmark must not only measure the performance of the storage system,
but \del{should} also provide a meaningful way to compare different systems and
configurations.
\add{For this reason, the HPC community maintains rankings in which every
system runs the same benchmarks under a common set of rules, so that the
resulting scores can be compared and ordered. The best known of them is the
TOP500~\cite{top500}, which ranks HPC platforms by their computational
performance. In the I/O world, the community created}
the IO500 benchmark suite\add{~\cite{io500list}.} \del{was designed for this
purpose.}
\add{IO500} is widely used to evaluate and rank the I/O performance of HPC storage
systems, and it combines several workloads, including bandwidth-oriented and
metadata-oriented operations, into a single score\replace{The benchmark results are also}{, which is} used to produce the IO500 ranking.\del{, which is
published twice a year.}

Because IO500 is a ranking, there is a direct incentive for HPC sites to
obtain the best possible score\add{. For vendors and system integrators, a high
position is a commercial argument, used in marketing material and in
procurement discussions. For computing centres, it is a public and comparable
demonstration that the machine they operate delivers the capability it was
built for, which matters both for the visibility of the centre and for
justifying the investment it represents}. \replace{
The}{However, the} reported performance of a system \del{therefore} does not only depend on its
hardware\replace{.
It also depends}{, but also} on how the system is configured and how the benchmark is
executed. \del{Finding a good configuration can consequently have a significant impact on the
measured IO500 score.} \add{Finding a good configuration is therefore essential
to obtain the highest possible score.}

\add{However, there is a} \del{the} large number of possible configurations\add{, which} makes the tuning of IO500
difficult\add{:} a configuration can involve several parameters, and each parameter can have
multiple possible values.
The combination of these parameters creates a potentially large search space,
and finding a good configuration therefore requires evaluating many possible
combinations.

Another important difficulty is the cost of an IO500 evaluation.
According to the IO500 submission rules\add{~\cite{io500rules}}, each write and
create phase must run
for at least \SI{300}{\second}\del{, using a stonewall timer set to 300
seconds}.
A run that finishes the write phase earlier is considered invalid.
This requirement makes a valid IO500 measurement relatively expensive in terms
of cluster time: on the target platform, one valid measurement takes
approximately \SI{30}{\minute} of exclusive cluster time.
When several configurations and repetitions are considered, the total cost
becomes very large.  \add{For instance,} in the parameter space considered in this work, approximately
\todoitem{confirm the exact number} $7{,}840$ configurations have to be
evaluated.
Measuring all of them with three repetitions would require approximately
\SI{12000}{\hour} of wall-clock time, which corresponds to more than one year
of computation on the target machine. \add{Therefore, } An exhaustive search is not practical.



\del{Instead, practitioners often rely on their experience and knowledge of the
system.
Another possible approach is to change one parameter at a time and select the
best value before moving to the next parameter.
Such approaches can work in some cases, but they may miss interactions between
parameters: a configuration that is optimal for one parameter independently is
not necessarily optimal when several parameters are considered together.}


This leads to the main research question of this internship:

\begin{quote}
    \itshape
    Can a smarter search method find a near-optimal IO500 configuration using
    far fewer evaluations, and which search method is most effective for this
    problem?
\end{quote}

\del{The goal is therefore not simply to find a configuration with a high IO500
score.
It is also important to understand how many evaluations are required to find
this configuration, and how different search strategies behave under a limited
evaluation budget.}

\subsection{Objectives of the internship}

\add{The objective of this internship is to tune IO500 with a guided search
instead of an exhaustive one, and to determine which search method should be
used in practice. Bayesian Optimization is the main candidate, but it is a
family of methods rather than a single one, so a large part of the work
consists in comparing its variants, and comparing them against random and
exhaustive search.}

\add{The work is carried out with IOPS~\cite{iops}, a benchmark orchestration
framework developed in the team. IOPS describes a parametric study in a single YAML file
and takes care of generating the runs, submitting them to the cluster
scheduler, parsing their output and caching the results. Two of its features
matter here: the search strategy is a configuration option, so exhaustive,
random and Bayesian searches can be compared without changing anything else;
and a study can be replayed from the cache, which is what makes the offline
comparison possible.}

\subsection{Contributions}

\add{The main contributions of this work are the following.}

\begin{itemize}
    \item \emph{A cheap ground-truth methodology} to evaluate search strategies
          for tuning IO500, based on a set of invalid scores that are much
          cheaper to generate but still allow the strategies to be compared.
    \item \emph{A systematic comparison of Bayesian Optimization variants},
          which identifies the setup that performs best and shows which of its
          parameters actually matter.
    \item \emph{A validation campaign} on the cluster, checking that the
          configurations selected on the cheap proxy remain the best ones
          under the official IO500 rules.
\end{itemize}

\subsection{Organisation of the report}

% Turn the section names into cross-references (Section~\ref{sec:...}) once the
% corresponding sections have been added to main.tex.
\add{The remainder of this report is organised as follows. The next section
presents the host laboratory and the environment in which the internship took
place. The background section then introduces what is needed to read the rest
of the report: the IO500 benchmark and its rules, the parameters exposed by
parallel file systems, and the principles of Bayesian Optimization. Related
work on I/O auto-tuning is reviewed next, followed by a description of the
tools, the platform and the experimental setup. The methodology section
explains how the ground truth was built and how the search strategies were
compared, and the results section reports and analyses the measurements. The
report closes with a discussion of the limitations of the study and a
conclusion.}
